\clearpage
\setlength{\parskip}{0.5em}

\markboth{\textbf{Abstract}}{\textbf{Abstract}}\label{abstract}

\textbf{Abstract}

\textit{Kann ein KI-gestützter Tutor den Einstieg ins Programmieren erleichtern?}


Die vorliegende Bachelorarbeit untersucht, inwiefern die Integration eines KI-gestützten Chat-Tutors und einer adaptiven Drill-Auswahl die Lernerfahrung beim Erwerb von Python-Grundlagen verbessert. Hierzu wurde ein webbasiertes Lernsystem entwickelt, das Einsteigende schrittweise durch die Implementierung eines vollständigen BMI-Rechner-Programms führt und dabei KI-basierte Assistenz über die OpenAI-API integriert. Als theoretischer Rahmen dient das 4C/ID-Modell (Four-Component Instructional Design) nach van Merriënboer, wobei der Fokus auf den Komponenten Learning Tasks und Part-task Practice liegt.

Eine kontrollierte empirische Studie mit 18 Teilnehmenden im Between-Subjects-Design untersuchte, ob KI-gestützte Komponenten zu günstigeren Lernindikatoren führen. Die Kontrollgruppe (Gruppe~A) arbeitete ohne KI-Funktionen, während die Experimentalgruppe (Gruppe~B) Zugang zu KI-gestützter Aufgabenauswahl und einem KI-Chat-Tutor erhielt. Die Datenerhebung umfasste System-Logs, NASA-TLX-Fragebögen, Usability-Skalen sowie KI-spezifische Evaluationsitems.

Die Ergebnisse zeigen keine signifikanten Unterschiede in objektiven Leistungsindikatoren wie Fehlerrate oder Abschlussquote. Auf subjektiver Ebene wies Gruppe~B konsistent günstigere Werte auf: geringere Frustration im NASA-TLX, höhere hedonische Qualität im UEQ-S sowie einen besseren Net Promoter Score. Eine qualitative Analyse der Chat-Verläufe identifizierte sowohl produktive Nutzungsmuster als auch systemseitige Schwachstellen, aus denen konkrete Verbesserungsmaßnahmen abgeleitet werden.

\setlength{\parskip}{0pt}
\clearpage
